%============================================================
% 01_abstract.tex — 국문 초록
%
% 초록은 논문 전체를 한 단락으로 요약한 글입니다.
% 연구 목적 · 방법 · 결과 · 결론의 흐름으로 작성하며,
% 분량은 보통 200~400자 내외입니다.
% 논문을 다 쓴 뒤 마지막에 작성하는 것을 권장합니다.
%============================================================

\begin{center}
  {\Large\textbf{국문 초록}}
  \noindent\rule{\textwidth}{.12mm}
\end{center}

초록은 언제 쓰는가? 맨 나중에 쓰는게 국룰이다. 진짜? 당연하지. 초록은 논문의 얼굴과 같다. 선행연구 \citep{smith1987tiling}에 의하면 어쩌고 저쩌고.
% ↑ 위 예시 문장을 지우고 실제 초록을 작성하세요.
% 본문에서 인용한 참고문헌은 초록에서도 \citep{키} 형식으로 표기할 수 있습니다.

\vskip5pt
\noindent{\bfseries 주제어:}\ \mykeywordsko
% 주제어는 00_info.tex의 \mykeywordsko에서 자동으로 가져옵니다.

\noindent\rule{\textwidth}{.12mm}
